\documentclass[a4paper,12pt]{article}
\usepackage[utf8]{inputenc}
\usepackage[english]{babel}
\usepackage{authblk}
\usepackage{graphicx}
\usepackage{mathptmx}
\usepackage[singlespacing]{setspace}
\usepackage[headheight=1in,margin=1in]{geometry}
\usepackage{fancyhdr}
\usepackage{tikz}
\usepackage{booktabs}
\usetikzlibrary{positioning, arrows.meta, fit, backgrounds, calc, shadows.blur}
\renewcommand{\headrulewidth}{0pt}
\pagestyle{fancy}
\bibliographystyle{unsrt}

\tikzset{
  datasource/.style={
    rounded corners, 
    draw=blue!40, 
    fill=blue!8,
    minimum width=2.2cm, 
    minimum height=0.8cm,
    font=\sffamily\small,
    blur shadow
  },
  process/.style={
    rounded corners=4pt,
    draw=teal!50,
    fill=teal!10,
    minimum width=3cm,
    minimum height=1cm,
    font=\sffamily\small\bfseries,
    blur shadow
  },
  output/.style={
    rounded corners,
    draw=red!50,
    fill=red!10,
    minimum width=3cm,
    minimum height=1cm,
    font=\sffamily\small\bfseries,
    blur shadow
  },
  arrow/.style={
    -{Stealth[length=3mm]},
    thick
  }
}

\makeatletter
\def\@maketitle{%
  \newpage

  \begin{center}%
  \let \footnote \thanks
    {\LARGE \@title \par}%
  \end{center}%
  \par
  \vskip 0.1em}
\makeatother

\chead{%
  CS2Italy – May 19-21, 2026 – Torino, Italy\\
  %
}

\graphicspath{{images/}}

\title{Automating Policy Classification: A Hybrid Machine Learning Approach for Living Databases}


\date{}

\begin{document}

\maketitle
\thispagestyle{fancy}

\begin{center}
\author{Nicola Cerutti \footnote{Bank of England}}
\end{center}

\section*{Introduction}

Social science research increasingly relies on large-scale policy databases to inform evidence-based decision-making \cite{zanti2021evidence, yang2024evidence}. However, these resources face a fundamental challenge: they are static snapshots of inherently dynamic data, which can quickly become outdated or exhibit coverage biases \cite{szomszor2022overton}. The conventional workflows often treat data curation as project-specific rather than as a generalizable methodological problem, yet classification tasks (assigning policies to categories, identifying objectives, extracting structured features from text) are fundamentally algorithmic and could be automated.

The Food System Policy Database (FSPD) \cite{lowder2025}, containing approximately 15,000 policies across 192 countries, is an example of this challenge: while providing insights on global food system governance, it requires continuous manual curation with domain expertise.

We present a hybrid computational approach that transforms the FSPD into a ``living database'' through automated machine learning (ML) based classification. Our two-stage architecture combines transformer-based language models for semantic text understanding with gradient boosting models for integrating structured metadata. Hybrid systems using transformer embeddings with engineered features have been shown to outperform embedding-only models \cite{tian2023best, levy2022comparison}. This design pattern can be extended to many social science datasets mixing text and tabular features, providing a template for attempting to establish sustainable data curation processes across computational social science.

\section*{Data and Preliminary Results}

\textbf{Data.} The FSPD aggregates policies from four source databases (FAPDA, PINE, GIfNA, Nourishing) spanning 195 countries. Each policy includes rich text (\textit{policydecision\_details}, average 603 characters) and structured metadata (country, geographic region, income group, year) alongside expert-assigned classifications: (1) \textit{segment} (11 categories: inputs, farming, trade, processing, etc.); (2) \textit{lever} (113 specific policy instruments, e.g. taxes, subsidies, bans, labelling, etc.); (3) \textit{type} (9 high-level intervention categories); and (4) three binary objectives (inclusion, environmental sustainability, dietary health). The dataset presents common challenges: severe class imbalance (246:1 ratio), multilingual text, and heterogeneous features mixing text with categorical and temporal metadata.
\newline

\noindent\textbf{Methodology.} We developed a two-stage hybrid architecture leveraging the complementary strengths of deep learning and gradient boosting (Figure \ref{fig:architecture}):

\textit{Stage 1: Semantic Embedding Extraction.} We fine-tune XLM-RoBERTa-base \cite{conneau2020}, a multilingual transformer model (278M parameters), using multi-task learning with six prediction heads sharing a common encoder. To address class imbalance, we implement inverse frequency weighting (class weights proportional to 1/frequency). Critically, we extract 768-dimensional contextual embeddings from the transformer's \texttt{[CLS]} token, capturing rich semantic representations.

\textit{Stage 2: Multimodal Classification with XGBoost.} The extracted embeddings are concatenated with structured metadata features (country, geographic region, income group, source database) and fed into XGBoost classifiers \cite{chen2016}, one per task. This hybrid design addresses a key limitation of transformer-only approaches: while language models excel at semantic understanding, they may not fully leverage structured categorical features that carry additional predictive signal. The complete system includes confidence-based quality control, with predictions below a confidence threshold (0.7) flagged for expert review.

\noindent
\begin{figure}[htbp]
\centering
\resizebox{\textwidth}{!}{%
\begin{tikzpicture}[node distance=1.2cm and 1.5cm]

% ---- DATA SOURCES ----
\node[datasource] (fapda) {FAPDA};
\node[datasource, right=of fapda] (pine) {PINE};
\node[datasource, right=of pine] (gifna) {GIfNA};
\node[datasource, right=of gifna] (dots) {...};
\node[datasource, right=of dots] (nour) {Nourishing};

% ---- PIPELINE (centered under GIfNA) ----
\node[process, below=1.8cm of gifna] (schema) {Schema Mapping};
\node[process, below=of schema] (embed) {XLM-RoBERTa Embeddings};
\node[process, below=of embed] (xgb) {XGBoost Classification};

% ---- OUTPUTS ----
\node[output, below left=1.8cm and 1.2cm of xgb] (web) {Living FSPD Web Interface};
\node[output, below right=1.8cm and 1.2cm of xgb] (mend) {Mendeley Data Repository};

% ---- ARROWS (with small caps labels) ----
\draw[arrow] (fapda.south) -- node[pos=0.55, left, font=\scriptsize\scshape] {} (schema.north);
\draw[arrow] (pine.south)  -- node[pos=0.55, left, font=\scriptsize\scshape] {} (schema.north);
\draw[arrow] (gifna.south) -- node[pos=0.45, right, font=\footnotesize\scshape] {extract} (schema.north);
\draw[arrow] (dots.south)  -- node[pos=0.55, right, font=\scriptsize\scshape] {} (schema.north);
\draw[arrow] (nour.south)  -- node[pos=0.55, right, font=\scriptsize\scshape] {} (schema.north);

\draw[arrow] (schema) -- node[right, font=\footnotesize\scshape] {transform} (embed);
\draw[arrow] (embed)  -- node[right, font=\footnotesize\scshape] {encode} (xgb);

\draw[arrow] (xgb) -- node[pos=0.45, left, font=\footnotesize\scshape] {publish} (web);
\draw[arrow] (xgb) -- node[pos=0.45, right, font=\footnotesize\scshape] {archive} (mend);

% ---- PANELS (background) ----
\begin{scope}[on background layer]
  % Data Sources panel
  \coordinate (srcNW) at ($(fapda.north west)+(-8pt, 22pt)$);
  \coordinate (srcSE) at ($(nour.south east)+(8pt,-10pt)$);
  \path[fill=blue!6, draw=black!15, rounded corners=8pt, line width=0.8pt]
    (srcNW) rectangle (srcSE);

  % Pipeline panel (wider to span more horizontally)
  \coordinate (pipeNW) at ($(schema.north west)+(-65pt, 22pt)$);
  \coordinate (pipeSE) at ($(xgb.south east)+(65pt,-10pt)$);
  \path[fill=teal!6, draw=black!15, rounded corners=8pt, line width=0.8pt]
    (pipeNW) rectangle (pipeSE);

  % Outputs panel
  \coordinate (outNW) at ($(web.north west)+(-20pt, 22pt)$);
  \coordinate (outSE) at ($(mend.south east)+(10pt,-10pt)$);
  \path[fill=red!6, draw=black!15, rounded corners=8pt, line width=0.8pt]
    (outNW) rectangle (outSE);
\end{scope}

% ---- PANEL TITLES ----
\node[font=\sffamily\bfseries\small, anchor=north]
  at ($(srcNW)+(6pt,-6pt)$) {Data Sources};

\node[font=\sffamily\bfseries\small, anchor=north]
  at ($(pipeNW)+(6pt,-6pt)$) {Pipeline};

\node[font=\sffamily\bfseries\small, anchor=north]
  at ($(outNW)+(6pt,-6pt)$) {Outputs};

\end{tikzpicture}
}
\caption{Two-stage hybrid architecture for automated policy classification. Policy data from multiple sources (FAPDA, PINE, GIfNA, Nourishing) undergoes schema mapping, is encoded via XLM-RoBERTa, classified via XGBoost (segment, lever, type, objectives). It is then published to the living database and data repository with confidence scores.}
\label{fig:architecture}
\end{figure}

\noindent\textbf{Results.} The multi-task transformer (Stage 1) achieved strong baseline performance on text-only classification. Multilingual evaluation confirmed robust performance across Spanish, French, and English texts without requiring translation. The hybrid Stage 2 system consistently improved performance across five of six tasks (Table \ref{tab:results}). Feature importance analysis revealed that approximately 99\% of predictive gain comes from the transformer embeddings, with structured metadata contributing marginally---suggesting that the policy text already encodes much of the geographic and institutional context. The largest gains were in \textit{type} classification (+3.4 F1 points) and binary \textit{inclusion} (+6.0 F1 points).

\begin{table}[h]
    \centering
    \small
    \begin{tabular}{l|cc|cc|cc}
        \toprule
         & \multicolumn{2}{c|}{\textbf{Stage 1} (text only)} & \multicolumn{2}{c|}{\textbf{Stage 2} (hybrid)} & \multicolumn{2}{c}{\textbf{$\Delta$}} \\
         Task & Acc & F1 & Acc & F1 & $\Delta$Acc & $\Delta$F1 \\
         \midrule
         \multicolumn{7}{c}{\textit{Multi-class Classification}} \\
         \midrule
         Segment & 83.3 & 83.9 & 86.5 & 86.5 & +3.3 & +2.6 \\
         Lever    & 63.6 & 63.7 & 66.0 & 66.2 & +2.4 & +2.5 \\
         Type     & 86.4 & 87.0 & 90.4 & 90.4 & +4.1 & +3.4 \\
         \midrule
         \multicolumn{7}{c}{\textit{Binary Objectives}} \\
         \midrule
         Inclusion          & 85.8 & 65.4 & 87.1 & 71.4 & +1.4 & +6.0 \\
         Environmental Tgt. & 95.9 & 82.8 & 95.7 & 82.4 & $-$0.3 & $-$0.4 \\
         Dietary Health     & 84.3 & 70.9 & 84.9 & 73.8 & +0.6 & +2.8 \\
         \bottomrule
    \end{tabular}
    \caption{Stage 1 (XLM-RoBERTa, text only) vs.\ Stage 2 (hybrid embeddings + metadata) performance (\%). Lower F1 scores for binary objectives reflect class imbalance. Stage 2 improves on 5 of 6 tasks.}
    \label{tab:results}
\end{table}

\section*{Discussion and Next Steps}

This work identifies a pathway for transforming expert-curated databases from static snapshots into (mostly) automatically maintained resources. The two-stage hybrid architecture is deliberately domain-agnostic: while applied here to food policy, the same pattern can be adapted to any social science corpus combining rich text with structured metadata. \newline

\noindent\textbf{Next Steps.} The classification pipeline is complete. Remaining work focuses on deployment: building the ingestion layer to fetch new policies from source databases (FAPDA, PINE, GIfNA), schema mapping to the FSPD format, deduplication, and a REST API with monthly automated updates via GitHub Actions. Initial predictions will be routed to domain experts for validation before full automation. Cross-database generalization tests will confirm that models learn genuine policy semantics rather than source-specific artifacts.

From a food policy perspective, a living FSPD will enable new longitudinal analyses: tracking policy diffusion across countries, measuring responses to shocks (pandemics, climate events, price crises), and identifying emerging trends in near-real-time. The original FSPD, for example, revealed that most food policies target producers and consumers while under-regulating the middle of the value chain \cite{lowder2025}; automated updates might reveal whether this pattern shifts as governments implement new policies.

By releasing the classification pipeline as open-source tooling, we aim at providing the scientific community with reusable infrastructure for automating data curation workflows, reducing barriers to longitudinal and real-time analysis of policy, legislative, and governance data across domains.

\bibliography{refs}


\end{document}

